% titel: Blatt 9, A1: Glukose/Ethanol/Wasser-Mischung
% Teil II: der durchgerechnete Weg.

\subsection*{Gegeben und gesucht}
20.0 g Glukose (C$_6$H$_{12}$O$_6$), 2.0 cm$^3$ Ethanol
($\rho=0.800\ g\,cm^{-3}$), 120.0 g H$_2$O; Dichte der Mischung
$\rho=0.940\ g\,cm^{-3}$. Gesucht: c, b und x der Komponenten.

\subsection*{Rechnung}
Molare Massen: $M(\text{Gluk}){=}180.16$, $M(\text{EtOH}){=}46.07$,
$M(\text{H}_2\text{O}){=}18.02$ (in g/mol).
\begin{align*}
n_{\text{Gluk}} &= 20.0/180.16 = 0.1110\ mol,\\
m_{\text{EtOH}} &= 0.800\cdot2.0 = 1.6\ g,
  \quad n_{\text{EtOH}} = 1.6/46.07 = 0.0347\ mol,\\
n_{\text{H}_2\text{O}} &= 120.0/18.02 = 6.66\ mol.
\end{align*}
Gesamtmasse $m=20.0+1.6+120.0 = 141.6$ g;
\[
  V = \frac{m}{\rho} = \frac{141.6}{0.940} \approx 150.6\ cm^3
    = 0.1506\ L.
\]

\subsubsection*{a) Stoffmengenkonzentrationen}
\[
  c_{\text{Gluk}} = \frac{0.1110}{0.1506} \approx 0.737\ mol\,L^{-1},
  \qquad
  c_{\text{EtOH}} = \frac{0.0347}{0.1506} \approx 0.231\ mol\,L^{-1}.
\]

\subsubsection*{b) Molalitäten} (Lösungsmittel = 0.120 kg Wasser)
\[
  b_{\text{Gluk}} = \frac{0.1110}{0.120} \approx 0.925\ mol\,kg^{-1},
  \qquad
  b_{\text{EtOH}} = \frac{0.0347}{0.120} \approx 0.289\ mol\,kg^{-1}.
\]

\subsubsection*{c) Molenbrüche}
$n_{\text{ges}} = 0.1110+0.0347+6.66 = 6.806$ mol.
\begin{align*}
x_{\text{Gluk}} &= 0.1110/6.806 \approx 0.0163,\\
x_{\text{EtOH}} &= 0.0347/6.806 \approx 0.0051,\\
x_{\text{H}_2\text{O}} &= 6.66/6.806 \approx 0.9786.
\end{align*}

\begin{ergebnis}
a) $c_{\text{Gluk}}\approx0.737$, $c_{\text{EtOH}}\approx0.231\ mol\,L^{-1}$.
b) $b_{\text{Gluk}}\approx0.925$, $b_{\text{EtOH}}\approx0.289\ mol\,kg^{-1}$.
c) $x_{\text{Gluk}}\approx0.016$, $x_{\text{EtOH}}\approx0.005$,
$x_{\text{H}_2\text{O}}\approx0.979$.
\end{ergebnis}

\begin{kontrolle}
Der Molebruch von Wasser dominiert (fast 98\,\%) --- erwartbar bei viel
Wasser. Summe $x=0.0163+0.0051+0.9786=1.000$: Probe erfüllt.
\end{kontrolle}

\begin{merksatz}
c (pro Liter Lösung), b (pro kg Lösungsmittel), x (pro Mol gesamt) ---
Nenner genau unterscheiden und die Lösung von $\rho$ holen.
\end{merksatz}
